% !Mode:: "TeX:UTF-8"
% !Mode:: "TeX:UTF-8"
\newcommand{\thepage}{} % to silence a hyperref warning about pdfpagelabels
\RequirePackage[]{silence}
\WarningsOff[hyperref]
\PassOptionsToPackage{log-declarations=false}{xparse}
%\RequirePackage[l2tabu, orthodox]{nag}
\RequirePackage{luatex85} % to make \pdf...= commands work with LuaTeX
\documentclass[aspectratio=169,t]{beamer}
\usepackage{ifthen}
\usepackage[T1]{fontenc}
\usepackage{iftex}
\ifLuaTeX
 \usepackage{fontspec}
 \setmainfont[Ligatures=TeX]{Source Sans Pro}
\else
 \usepackage[utf8]{inputenc}
\fi
\usepackage{german}
\usepackage{newunicodechar}
\newunicodechar{€}{\texteuro}
\newunicodechar{ß}{\ss}
\renewcommand{\SS}{\iffontchar\font"1E9E \symbol{"1E9E}\else SS\fi}
\usepackage{ellipsis}
\usepackage{microtype}
\usepackage{listings}
\usepackage{sourcecodepro}
\lstset{language=Python, basicstyle=\ttfamily}
\usepackage{amsfonts}
\usepackage{amssymb}
\usepackage{amsthm}
\usepackage{fancyhdr}
\usepackage{graphicx}

\newtheorem{proposition}[theorem]{Proposition}

\DeclareMathOperator*{\argmax}{argmax}
\DeclareMathOperator{\supp}{supp}
\newcommand{\E}{\ensuremath{\mathbb{E}}}
\newcommand{\V}{\ensuremath{\mathbb{V}}}

\setbeamercolor{emph}{fg=purple}
\renewcommand<>{\emph}[1]{%
  {\usebeamercolor[fg]{emph}\only#2{\itshape}#1}%
}



\hypersetup{draft=false,colorlinks=false,pdfpagemode=UseOutlines,plainpages=false,
            breaklinks=true,bookmarksnumbered=true,pdfborder=0 0 0,
            pdftitle={Data Science SoSe 2026 HHU},
            pdfauthor={Konrad Voelkel},
            pdfsubject={data science}}

%\documentclass[aspectratio=169,t]{beamer}
\usetheme{hhubeamer} % put the hhu-beamer-theme in /usr/share/texmf/tex/latex/beamer/themes/

% -------------------- Code listings config --------------------
% Define custom colors for syntax highlighting
\definecolor{codegreen}{rgb}{0,0.6,0}
\definecolor{codegray}{rgb}{0.5,0.5,0.5}
\definecolor{codepurple}{rgb}{0.58,0,0.82}
\definecolor{backcolour}{rgb}{0.96,0.96,0.94} % Very light gray/beige
% Configure the listings package
\lstset{
    language=Python,
    backgroundcolor=\color{backcolour},   
    commentstyle=\color{codegreen}\itshape,
    keywordstyle=\color{blue}\bfseries,
    numberstyle=\tiny\color{codegray},
    stringstyle=\color{codepurple},
    basicstyle=\ttfamily\footnotesize, % Use a monospaced font
    breakatwhitespace=false,         
    breaklines=true,      
             % Enable automatic line breaking
    captionpos=b,                      % Caption at the bottom
    keepspaces=true,                 
    numbers=left,                      % Line numbers on the left
    numbersep=10pt,                    % Distance of line numbers from the code
    showspaces=false,                
    showstringspaces=false,            % Don't show weird underscores in strings
    showtabs=false,                  
    tabsize=4,
    frame=single,                      % Add a subtle border around the code
    rulecolor=\color{black!20},        % Border color
  }

\event[]{}
\author[Sommersemester 2026]{Dr. Konrad Völkel, Sommersemester 2026}

%\title{Data Science, VL 1}
%\date{16.04.2026}


%%% Local Variables:
%%% mode: latex
%%% TeX-master: t
%%% End:


\title{Data Science, VL 4} % war in 2025 noch VL 3; Einschub zu Python ...
\date{28.04.2026} %\date{23.04.2026}

\begin{document}

\begin{frame}
  \titlepage
\end{frame}

\section{Recap}

\begin{frame}
  \frametitle{Was bisher geschah: Stochastik}
  \begin{itemize}[<+->]
  \item Ein \alert{Wahrscheinlichkeitsmaß} auf einer Menge $\Omega$ ist Abbildung $P : \mathcal{E}(\Omega) \to [0,1]$ mit $P(\Omega)=1$ und Summenregel für disjunkte Vereinigungen. Wir nennen $A \in \mathcal{E}(\Omega)$ die \alert{Ereignisse}, für eine endliche Menge $\mathcal{E}(\Omega) := \mathcal{P}(\Omega)$, die Menge aller Teilmengen. Das Paar $(\Omega,P)$ heißt \alert{W-Raum}.
    \item Auf einer endlichen Menge $\Omega$ ist $P$ durch eine \alert{pmf} $p \colon \Omega \to [0,1]$ gegeben, $P(A) = \sum_{x\in A} p(x)$.
    \item \alert{Zufallsvariablen} sind messbare Abbildungen $X \colon \Omega' \to \Omega$, die induzierte Verteilung $P_X$ heißt auch \alert{Verteilung von $X$}.
    \item Sind $\Omega, \Omega'$ endlich, so sind alle Abbildungen $X \colon \Omega' \to \Omega$ messbar.
    \end{itemize}
\end{frame}



\section{Beispiel Daten, Stichprobe, Modell}

\begin{frame}
  \frametitle{Daten als Verteilung}
 \uncover<+->{
   Gegeben Daten $D$ in einer Tabelle:
   
   \begin{tabular}{|c|c|}
     \hline
       & x \\ \hline
     0 & blue \\
     1 & red \\
     2 & blue \\
     3 & green \\ \hline
   \end{tabular}
 }

 \uncover<+->{
   können wir einen W-Raum $\Omega = \{\text{blue}, \text{red}, \text{green}\}$ modellieren mit
   \[P(\{x\}) := \text{relative Frequenz von } x = \dfrac{\text{Anzahl } x \text{ in } D }{|D|}\]
   Die \alert{Verteilung der Daten} (und $|D|=4$ im Beispiel).
 }
 
\end{frame}

\begin{frame}
  \frametitle{Daten als Verteilung}
   Gegeben Daten $D$ in einer Tabelle:
   
   \begin{tabular}{|c|c|}
     \hline
       & x \\ \hline
     0 & blue \\
     1 & red \\
     2 & blue \\
     3 & green \\ \hline
   \end{tabular}

   ist der Wahrscheinlichkeitsvektor für die Reihenfolge $\Omega = \{\text{blue}, \text{red}, \text{green}\}$
   \[p = \left(\frac{2}{4}, \frac{1}{4}, \frac{1}{4}\right).\]
 
\end{frame}

\begin{frame}
  \frametitle{Populationsverteilung}
  \begin{itemize}[<+->]
    \item  Wo kommen die Daten her?
 \item  Beispiel: Stichprobe einer Tüte Süssigkeiten. \\
  Achtung: größere Stichproben können anders verteilt sein.
\item Häufige Aufgabe: Gegeben Stichprobe $D$ einer größeren Menge $D'$,
  schätze Verteilung von $D'$ anhand der Verteilung der Daten $D$.
\item  Dabei heißt $D'$ die \alert{Population}.
\item  Häufig zusätzlicher Input: \\
  Zusatzinformationen über die Verteilung von $D'$ (eine Art von \alert{Prior}).
\end{itemize}
\end{frame}

\begin{frame}
  \frametitle{Modellverteilung}
  \begin{itemize}[<+->]
  \item Wenn z.B. die Modellannahme ist, dass die Farben gleichverteilt sind bis auf eine Farbe, die doppelt so oft vorkommt,
    dann sind wir geneigt, nicht die Verteilung der Daten als Modell zu wählen, sondern z.B.
    \[p = \left(\frac{2}{4}, \frac{1}{4}, \frac{1}{4}\right),\]
    auch wenn die konkreten relativen Häufigkeiten leicht davon abweichen.
  \item Auch mit dem Modell lässt sich eine ``Stichprobe'' ziehen: \\
    synthetische Daten, mit Hilfe eines Zufallszahlgenerators.
  \item Übungsaufgabe: wie genau geht das im Beispiel?
\end{itemize}
\end{frame}

\begin{frame}
  \frametitle{Modellgüte}
  Wie gut ist das Modell?
  
  Kann vergleichen:\\
  sehen die Daten $D$ so aus, als wären sie aus dem Modell gesampelt worden?
  
  \uncover<2->{``How likely is that?'' $\to$ \alert{Likelihood}.}
\end{frame}


\begin{frame}
  \frametitle{Was bisher geschah: Statistik}
  \begin{itemize}[<+->]
  \item Wir wählen einen Ausschnitt der Welt, sammeln Daten, putzen, transformieren (d.h. wählen die \alert{Features} aus, die uns interessieren) und bilden Modelle für die Daten: Wahrscheinlichkeitsverteilungen.
  \item \alert{Maximum Likelihood-Prinzip}: Wir wählen das Modell, zu welchem die beobachteten Daten am ehesten so aussehen, als hätten wir sie aus dem Modell synthetisch erzeugt (\alert{gesampelt}).
    \end{itemize}
    \includegraphics[width=12cm]{figures/dsprozess-simplified-2.png}
\end{frame}




\begin{frame}[separator]
  \frametitle{Pause?}
\end{frame}



\section{Abweichungen und Mittel}


\subsection{Herleitung}

\begin{frame}
  \frametitle{Vergleich von Verteilungen}
  \begin{itemize}[<+->]
  \item Wie messen wir den Unterschied der Verteilungen aus den Daten und der Verteilungen aus unseren Modellen?
  \item Präzise Antwort: mit der relativen Entropie - das kommt später.
  \item Möglichkeit: Abstand der Daten zu synthetischen Daten im Datenraum messen.
    \item Damit wir wissen, wie der Datenraum aussieht: Ab jetzt $X$ eine reelle Zufallsvariable bzw. Daten $x$ sind reelle Zahlen.
    \end{itemize}
  \end{frame}

  \begin{frame}
    \frametitle{Vergleich von Daten}
    Gegeben $(x,y)$ mit Stichprobengröße $n=4$ in zwei Spalten
    
    (Idee: $x$ sind die Daten, $y$ unser Modell, unsere ``Vorhersage'')
    
    \begin{tabular}{|c|c|c|}
     \hline
       & x & y \\ \hline
     0 & 178 & 178 \\
     1 & 160 & 155 \\
     2 & 155 & 150 \\
     3 & 174 & 180 \\ \hline
    \end{tabular}
    
    können wir nach dem Abstand von $x$ und $y$ fragen:
    \[d(x,y) = \Vert x - y \Vert_1 = \sum_{i=0}^{n-1} |x_i-y_i| \]
  \end{frame}

    \begin{frame}
    \frametitle{Vergleich von Daten}
\begin{itemize}[<+->]
\item    Gegeben Daten $(x,y)$ mit Stichprobengröße $n=4$ können wir nach dem Abstand von $x$ und $y$ fragen:
  \[d(x,y) = \Vert x - y \Vert_1 = \sum_{i=0}^{3} |x_i-y_i| \]
  \uncover<1>{Achtung: Das wächst mit zunehmendem $n$}
\item    Ebenso möglich wäre der \alert{mittlere Fehler}:
  \[\dfrac{1}{n} \sum_{i=0}^{n-1} (x_i-y_i)\]
\end{itemize}
\end{frame}

    \begin{frame}
    \frametitle{Vergleich von Daten}
\begin{itemize}[<+->]
\item    Gegeben Daten $(x,y)$ mit Stichprobengröße $n$ ist der Abstand bzgl. $1$-Norm:
  \[d_1(x,y) = {\Vert x - y \Vert}_1 = \sum_{i=0}^{n-1} |x_i-y_i| \]
  Der \alert{mittlere absolute Fehler}
  $mae(x,y) = \dfrac{d_1(x,y)}{n} $
\item Der Abstand bzgl. $2$-Norm ist (quadriert)
  \[d_2(x,y)^2 = {\Vert x - y \Vert}_2^2 = \sum_{i=0}^{n-1} {\left(x_i-y_i\right)}^2 \]
  Der \alert{mittlere quadratische Fehler}
 $mse(x,y) = \dfrac{d_2(x,y)^2}{n} $
\end{itemize}
\end{frame}

  \begin{frame}
    \frametitle{Beispielvergleich von Daten}
    Gegeben $(x,y)$ mit Stichprobengröße $n=4$ in zwei Spalten
    
    (Idee: $x$ sind die Daten, $y$ unser Modell, unsere ``Vorhersage'')
    
    \begin{tabular}{|c|c|c|c|}
     \hline
       & x & y & x - y\\ \hline
     0 & 178 & 178 & 0 \\
     1 & 160 & 155 & 5\\
     2 & 155 & 150 & 5\\
     3 & 174 & 180 & -6\\ \hline
    \end{tabular}

    \[d_1(x,y) = 0+5+5+6 = 16,\ d_2(x,y)^2 = 0^2+5^2+5^2+6^2 = 86\]
    \[mae(x,y) = \dfrac{16}{4} = 4,\ mse(x,y) = \dfrac{86}{4} = 21.5\]
    Der mittlere Fehler ist $\dfrac{0+5+5-6}{4} = 1$, eine trügerisch geringe Zahl.
  \end{frame}

  \begin{frame}
    \frametitle{Beispielvergleich von Daten: einfacher Fall}
    Wir wählen nun ein besonders einfaches Modell:

    ein konstanter Wert (hier $165$) als Vorhersage.
    
    \begin{tabular}{|c|c|c|c|}
     \hline
       & x & y & x - y\\ \hline
     0 & 178 & 165 & 13 \\
     1 & 160 & 165 & -5\\
     2 & 155 & 165 & -10\\
     3 & 174 & 165 & 9\\ \hline
    \end{tabular} \uncover<2->{Für welche Konstante wird der Fehler minimal?}

    \[d_1(x,y) = 13+5+10+9 = 37,\ d_2(x,y)^2 = 13^2+5^2+10^2+9^2 = 375\]
    \[mae(x,y) = \dfrac{37}{4} = 9.25,\ mse(x,y) = \dfrac{375}{4} = 93.75\]
    Der mittlere Fehler ist $\dfrac{13-5-10+9}{4} = 7$, eine trügerisch geringe Zahl.
  \end{frame}

\subsection{Mittelwerte}
  
    \begin{frame}
    \frametitle{Vergleich von Daten mit einer Konstanten}
\begin{itemize}[<+->]
\item    Gegeben Daten $x$ mit Stichprobengröße $n$ und einer Konstanten $y$ können wir nach dem Fehler fragen, den wir machen, wenn wir $x$ durch $y$ modellieren.
\item  Ein $y$, welches den mittleren absoluten Fehler $mae(x,y)$ minimiert, nennt man \alert{Median} von $x$,
  $median(x) \in argmin_{c \in \mathbb{R}} mae(x,c)$. Manche fordern $c \in x$, dass also der Median ein Datenpunkt sein muss.
\item Ein $y$, welches den mittleren quadratischen Fehler $mse(x,y)$ minimiert, nennt man \alert{Mittelwert} von $x$, $mean(x) = argmin_{c \in \mathbb{R}} mse(x,c)$
\item Es ist $mean(x) = \frac{1}{n} \sum_{i=0}^{n-1} x_i$, das arithmetische Mittel.
  \item Für ungerades $n$ schreiben wir $n=2m+1$. Wenn $i(0),\dots,i(2m)$ Indizes sind sodass $x_{i(j)} \leq x_{i(j+1)}$, dann ist das mittlere $x_{i(m)}$ ein Median von $x$.
\end{itemize}
\end{frame}

  \begin{frame}
    \frametitle{Beispielvergleich von Daten mit konstantem Modell}
    Wir wählen nun konstante Modelle:
    Median $167$ und Mittelwert $166.75$.
    
    \begin{tabular}{|c|c|c|c|}
     \hline
       & x & x - median(x) & x - mean(x)\\ \hline
     0 & 178 & 11 & 11.25 \\
     1 & 160 & -7 & -6.75\\
     2 & 155 & -12 & -11.75\\
     3 & 174 & 7 & 7.25\\ \hline
    \end{tabular}
    \[mae(x,median(x)) = \dfrac{11+7+12+7}{4} = 9.25,\ mse(x,mean(x)) = 90.6875 \approx 9.523^2\]
    Es gibt keine konstanten Modelle mit kleineren $mae$ bzw. $mse$.
  \end{frame}

\subsection{Streumaße}
  
   \begin{frame}
    \frametitle{Abweichung vom besten konstanten Modell}
\begin{itemize}[<+->]
\item Die Größe $mse(x,mean(x))$ ist ein Maß für die Streuung der Daten.
\item Da der mittlere quadratische Fehler auf einer anderen Skala liegt als die Daten selbst, ziehen wir die Wurzel:
  \[\sigma(x) := rmse(x,mean(x)) := \sqrt{mse(x,mean(x))}\]
  heißt \alert{root mean squared error} oder auch \alert{Stichproben-Standardabweichung}.
\item Nützliche Formel (wobei $x^2$ komponentenweise zu verstehen ist):
  \[\sigma(x)^2 = \frac{1}{n} \sum_{i=0}^{n-1} {\left( x_i - mean(x) \right)}^2 = mean(x^2)-mean(x)^2\]
\end{itemize}
\end{frame}

  \begin{frame}
  \frametitle{Zusammenfassung: Abweichungen und Mittel}
  \begin{itemize}
    \item \textbf{Ziel:} Unterschied zwischen Daten und Modell quantifizieren.
    \item \textbf{Abstandsmaße:}
      \begin{itemize}
        \item \textbf{Mittlerer absoluter Fehler (MAE):} mittlerer Betrag der Abweichungen
        \item \textbf{Mittlerer quadratischer Fehler (MSE):} mittleres Quadrat der Abweichungen
      \end{itemize}
    \item \textbf{Beste konstante Modelle:}
      \begin{itemize}
        \item \textbf{Median} minimiert den MAE
        \item \textbf{Mittelwert} minimiert den MSE
      \end{itemize}
    \item \textbf{Streumaß:}
      \begin{itemize}
        \item \textbf{Standardabweichung (RMSE):} Wurzel des MSE beim Mittelwert-Modell
      \end{itemize}
  \end{itemize}
\end{frame}


\begin{frame}[separator]
  \frametitle{Pause?}
  
  \vspace{-2.6cm}
\includegraphics[height=\textheight]{figures/gambling-dogs.png}
\end{frame}


\section{Momente einer Zufallsvariablen}

\subsection{Erwartungswert}

   \begin{frame}
    \frametitle{Mittelwerterwartung}
\begin{itemize}[<+->]
\item Das konstante Modell ist eine Dirac-Verteilung. Für eine interessantere Verteilung können wir Stichproben ziehen (sampeln) und deren Mittelwerte und Standardabweichungen berechnen. Der Vergleich mit Mittelwert und Standardabweichung der Daten gibt Anhaltspunkte für gute Modelle.
\item Je größer die synthetischen Daten, desto aussagekräftiger.
\item Es geht noch besser: wir berechnen den Grenzwert des ``synthetischen Mittelwerts'' ohne je eine Stichprobe simulieren zu müssen. Wie?
\end{itemize}
\end{frame}

   \begin{frame}
    \frametitle{Mittelwerterwartung}
    \begin{itemize}[<+->]
    \item Wenn $p = \left( \frac{a_1}{N}, ..., \frac{a_k}{N} \right)$ ein Wahrscheinlichkeitsvektor ist (also $k$ mögliche Merkmalsausprägungen $y_1,\dots,y_k \in \mathbb{R}$),
      so hätte ein ``idealisiertes Sample'' $\tilde{D}$ z.B. die Häufigkeiten $(a_1,\dots,a_k)$.
      \item Der Mittelwert wäre dann $mean(\tilde{D}) = \frac{1}{N} \sum_{j=1}^k a_j y_j$.
      \item Allgemein definieren wir für eine reelle Zufallsvariable $X \colon \Omega \to \mathbb{R}$ mit nur endlich vielen oder höchstens abzählbar unendlich vielen Werten
        \[\mathbb{E}(X) := \sum_{\omega \in \Omega} P(\{\omega\}) X(\omega) = \sum_{x \in \mathbb{R}} p_X(x)x\]
        ($p_X$ eine pmf für die Verteilung von $X$).
        Das ist der \alert{Erwartungswert} von $X$.
\end{itemize}
\end{frame}

   \begin{frame}
    \frametitle{Erwartungswert einer Gleichverteilung}
    \begin{itemize}[<+->]
    \item Sei $\Omega = \{0,\dots,n-1\}$ mit einer Gleichverteilung gegeben und $X \colon \Omega \to \mathbb{R}$ die Abbildung $\omega \mapsto \omega+1$. Dann hat die Verteilung von $X$ den Erwartungswert
      \[\mathbb{E}(X) = \sum_{\omega=0}^{n-1} \frac{\omega+1}{n} = \sum_{i=1}^{n} \frac{i}{n}  = \frac{n(n+1)}{2n} = \frac{n+1}{2}\]
      \item Wenn stattdessen $Y \colon \Omega \to \mathbb{R}$ mit $\omega \mapsto \omega$ eine Gleichverteilung auf $\{0,\dots,n-1\}$ beschreibt, so ist $\mathbb{E}(Y) = \frac{n+1}{2}-1$.
\end{itemize}
\end{frame}

\subsection{Varianz}

   \begin{frame}
    \frametitle{Streuungserwartung}
\begin{itemize}[<+->]
\item Das konstante Modell ist eine Dirac-Verteilung, die Streuung ist $0$. Für eine interessantere Verteilung können wir Stichproben ziehen (sampeln) und deren Standardabweichungen berechnen.
\item Wir berechnen den Grenzwert der ``synthetischen Standardabweichung'' ohne je eine Stichprobe simulieren zu müssen. Wie?
\item $\mathbb{V}(X) := \mathbb{E}\left({\left(X-\mathbb{E}X\right)}^2\right) = \mathbb{E}(X^2) - {\left(\mathbb{E}X\right)}^2$ heißt \alert{Varianz} von $X$
  und\\
  $\sigma(X) := \sqrt{\mathbb{V}X}$ die \alert{Standardabweichung} von $X$.
\item Tschebyscheff-Ungleichung: für $\lambda > 0$ ist
  $P\left(|X - \mathbb{E}(X)| \geq \lambda \sigma(X)\right) \leq \frac{1}{\lambda^2}$,
  die Wahrscheinlichkeit einer absoluten Abweichung um mehr als $\lambda$ Standardabweichungen liegt also unter $\frac{1}{\lambda^2}$.
\end{itemize}
\end{frame}

\subsection{Zusammenfassung}

\begin{frame}
  \frametitle{Zusammenfassung: Momente einer Zufallsvariablen}
  \begin{itemize}
    \item \textbf{Erwartungswert:}
      \begin{itemize}
        \item Mittelwert einer Zufallsvariablen: \(\mathbb{E}(X) = \sum_x p_X(x)\,x\)
        \item Der typische Wert, den man im Mittel bei vielen Stichproben erwartet
      \end{itemize}
    \item \textbf{Varianz und Standardabweichung:}
      \begin{itemize}
        \item Streuung um den Erwartungswert: \(\mathbb{V}(X) = \mathbb{E}(X^2) - \mathbb{E}(X)^2\)
        \item \(\sigma(X) := \sqrt{\mathbb{V}(X)}\) ist die Standardabweichung
      \end{itemize}
    \item \textbf{Interpretation:}
      \begin{itemize}
        \item $\mathbb{E}X$ und $\sigma(X)$ geben zentrale Lage und Streuung der Verteilung an
        \item \textbf{Tschebyscheff-Ungleichung:}\\
          Abweichungen um mehr als \(\lambda\) Standardabweichungen sind selten
      \end{itemize}
  \end{itemize}
  \vspace{0.5em}
  Erwartungswert und Varianz sind fundamentale Kennzahlen von Modellen.
\end{frame}

%    \begin{frame}
%     \frametitle{Ausblick: Höhere Momente}
% \begin{itemize}[<+->]
% \item Für allgemeines $n$ ist $\mu_n(X) := \mathbb{E}\left( X^n \right)$ das \alert{$n$-te Moment},
%  \item $\tilde{X} := X - \mathbb{E}(X)$ die \alert{zentrierte Zufallsvariable},
%  \item  $\tilde{\mu}_n(X) := \mu_n(\tilde{X})$ das \alert{$n$-te zentrale Moment}.
%    \item Es ist $\mu_1(X)=\mathbb{E}X$, $\tilde{\mu}_1(X) = 0$, $\tilde{\mu}_2(X) = \mathbb{V}X$.
% \end{itemize}
% \end{frame}

\begin{frame}[separator]
  \frametitle{Vielen Dank für Ihre Aufmerksamkeit!}
\end{frame}

\end{document}

%%% Local Variables:
%%% mode: latex
%%% TeX-master: t
%%% End:
